% ===== PRÉAMBULE CANONIQUE — à coller VERBATIM en tête de CHAQUE .tex =====
\documentclass[11pt,a4paper]{article}
\usepackage[utf8]{inputenc}\usepackage[T1]{fontenc}\usepackage{lmodern}
\usepackage[french]{babel}
\usepackage{amsmath,amssymb,mathtools}\usepackage{icomma}
\usepackage{siunitx}\sisetup{locale=FR}  % \qty{}{}, \si{}, \ang{}
\usepackage{xcolor}\usepackage{colortbl}
\usepackage{tikz}\usepackage{pgfplots}\pgfplotsset{compat=1.18}
\usepackage[siunitx,europeanresistors]{circuitikz} % circuits (élec.)
\usepackage[most]{tcolorbox}\usepackage{enumitem}\usepackage{array,booktabs}
\usepackage[a4paper,margin=2.2cm]{geometry}
\usetikzlibrary{arrows.meta,calc,angles,quotes,positioning,patterns,%
  backgrounds,shapes.geometric,decorations.pathmorphing,decorations.markings,intersections}
\definecolor{primary}{RGB}{222,49,99}\definecolor{primarydark}{RGB}{150,22,55}
\definecolor{defcol}{RGB}{0,114,178}\definecolor{methcol}{RGB}{52,92,148}
\definecolor{excol}{RGB}{0,135,90}\definecolor{attncol}{RGB}{200,40,55}
\tcbset{boxbase/.style={enhanced,breakable,boxrule=0.9pt,arc=2.5pt,
  left=3mm,right=3mm,top=2.3mm,bottom=2.3mm,fonttitle=\bfseries,coltitle=white}}
% --- boîtes TUTORIEL (noms EXACTS, ne pas renommer) ---
\newtcolorbox{cle}[1][]{boxbase,colback=defcol!6,colframe=defcol,
  title={\textsc{À retenir}\ifx\relax#1\relax\relax\else\textmd{ : \itshape #1}\fi}}
\newtcolorbox{methode}[1][]{boxbase,colback=methcol!7,colframe=methcol,sharp corners,
  title={\textsc{Méthode}\ifx\relax#1\relax\relax\else\textmd{ --- \itshape #1}\fi}}
\newtcolorbox{exemplebox}[1][]{boxbase,colback=excol!5,colframe=excol,
  title={\textsc{Exemple}\ifx\relax#1\relax\relax\else\textmd{ : \itshape #1}\fi}}
\newtcolorbox{piege}[1][]{boxbase,colback=attncol!6,colframe=attncol,
  title={\textsc{Piège}\ifx\relax#1\relax\relax\else\textmd{ : \itshape #1}\fi}}
\newtcolorbox{remarque}[1][]{boxbase,colback=black!4,colframe=black!45,
  title={\textsc{Remarque}\ifx\relax#1\relax\relax\else\textmd{ : \itshape #1}\fi}}
% --- boîtes EXERCICE / CORRIGÉ (noms EXACTS, auto-comptées) ---
\newtcolorbox[auto counter]{exo}[1][]{boxbase,colback=primary!4,colframe=primary,
  title={\textsc{Exercice}~\thetcbcounter\ifx\relax#1\relax\relax\else\textmd{ --- \itshape #1}\fi}}
\newtcolorbox[auto counter]{corrige}[1][]{boxbase,colback=excol!4,colframe=excol,
  title={\textsc{Corrigé}~\thetcbcounter\ifx\relax#1\relax\relax\else\textmd{ --- \itshape #1}\fi}}
\newcommand{\vv}[1]{\vec{#1}}\newcommand{\ds}{\displaystyle}
\newcommand{\R}{\mathbb{R}}\newcommand{\N}{\mathbb{N}}\newcommand{\Z}{\mathbb{Z}}
% =========================================================================
\begin{document}
\sisetup{per-mode=symbol}
% ================= FACILE =================

\begin{exo}[la pesanteur au sol]
Calcule l'intensité de pesanteur au niveau du sol $g_0=\dfrac{\mathcal{G}M_T}{R_T^2}$,
avec $\mathcal{G}=\num{6.67e-11}$ (SI), $M_T=\qty{6.0e24}{\kilogram}$ et
$R_T=\qty{6.4e6}{\metre}$.
\end{exo}

\begin{exo}[distance au centre]
Le rayon terrestre vaut $R_T=\qty{6400}{\kilo\metre}$. Donne la distance $d=R_T+h$
entre l'objet et le centre de la Terre pour :
\begin{enumerate}[label=\alph*)]
  \item une altitude $h=\qty{600}{\kilo\metre}$ ;
  \item une altitude $h=\qty{1600}{\kilo\metre}$ ;
  \item l'altitude géostationnaire $h=\qty{35800}{\kilo\metre}$.
\end{enumerate}
\end{exo}

\begin{exo}[pesanteur en altitude]
Un satellite est à l'altitude $h=\qty{3600}{\kilo\metre}$. On donne
$R_T=\qty{6400}{\kilo\metre}$ et $\mathcal{G}M_T=\qty{4.0e14}{}$ (SI).
Calcule la pesanteur $g_h$ à cette altitude.
\end{exo}

\begin{exo}[poids sur la Terre et sur la Lune]
Un astronaute a une masse $m=\qty{60}{\kilogram}$. On prend
$g_0=\qty{9.8}{\metre\per\second\squared}$ (Terre) et
$g_{\text{Lune}}=\qty{1.6}{\metre\per\second\squared}$.
\begin{enumerate}[label=\alph*)]
  \item Calcule son poids sur la Terre.
  \item Calcule son poids sur la Lune.
  \item Quelle est sa masse sur la Lune ?
\end{enumerate}
\end{exo}

\begin{exo}[du poids à la masse]
Un objet pèse $P=\qty{500}{\newton}$ sur Terre, où l'on prend
$g=\qty{10}{\metre\per\second\squared}$.
\begin{enumerate}[label=\alph*)]
  \item Quelle est sa masse ?
  \item Combien pèse-t-il sur Mars, où $g_{\text{Mars}}=\qty{3.7}{\metre\per\second\squared}$ ?
  \item Quelle est sa masse sur Mars ?
\end{enumerate}
\end{exo}

\end{document}
